\documentclass[10pt,letterpaper]{article} % exam class defaults to 10pt on letterpaper
\usepackage[margin=1in]{geometry}         % exam class defaults to 1-inch margins
\usepackage{graphicx} 
\usepackage{fancyhdr}
\usepackage{amssymb}
\usepackage{amsmath}
\usepackage{mathtools}
\usepackage[most]{tcolorbox}
\usepackage[dvipsnames]{xcolor}
\usepackage{blindtext}
\usepackage{amsthm}
\usepackage{upgreek}
\usepackage{hyperref}
\usepackage{tikz}
\usepackage{tikz-cd}
\usetikzlibrary{positioning}
%making an edit, want to see if auto script commits it%
%making another edit to test the commit bot

\setlength{\parindent}{0pt}
\setlength{\parskip}{1ex plus 0.5ex minus 0.2ex} % Matches the paragraph spacing of exam
\hbadness=10000 % Place in your preamble


\newtheorem{thm}{Theorem}[section]
\newtheorem{lem}{Lemma}[section]
\newtheorem{defn}{Definition}[section]
\newtheorem{cor}{Corollary}[section]
\newtheorem{axm}{Axiom}[section]

\tcolorboxenvironment{thm}{colback=lime, colframe=black, breakable}
\tcolorboxenvironment{lem}{colback=yellow, colframe=black, breakable}
\tcolorboxenvironment{defn}{colback=pink, colframe=black, breakable}
\tcolorboxenvironment{cor}{colback=SpringGreen, colframe=black, breakable}
\tcolorboxenvironment{axm}{colback=white, colframe=black, breakable}

% --- Header/Footer Setup ---
\fancyhf{}
\renewcommand\headrulewidth{0.5pt}
\pagestyle{fancy}
\author{Omkar Tasgaonkar}
\setlength{\headsep}{0.5in}
\fancyfoot[C]{(\thepage)}
\renewcommand\footrulewidth{0.5pt}

\title{Linear Algebra}
\author{Omkar Tasgaonkar \\ \small otasgaonkar@ucla.edu}
\date{}

\fancypagestyle{firstpage}{
    \fancyhf{}
    \fancyhead[C]{\textbf{Notes in Algebra}}
    \renewcommand\headrulewidth{0.5pt}
    \fancyfoot[C]{(\thepage)}
    \renewcommand\footrulewidth{0.5pt}
}
\begin{document}
\maketitle
This is a description of the notes and intended goals/topics that will be covered in them.

\begin{tcolorbox}[colback=white, colframe=black, arc=0mm]
Textbooks Used to Develop these Notes:
\begin{enumerate}
\item \textit{\textbf{Linear Algebra} - Hoffman and Kunze}
\end{enumerate}
Other Textbooks Referenced:
\begin{enumerate}
\item \textit{\textbf{Textbook} - Author}
\end{enumerate}
\end{tcolorbox}
\begin{tcolorbox}[colback=white, colframe=black, arc=0mm]
\textbf{Notes Last Updated:} \today
\end{tcolorbox}
\thispagestyle{firstpage}
\newpage
\begin{tcolorbox}[colback=white, colframe=black, arc=0mm]
\tableofcontents
\end{tcolorbox}
\textbf{NOTE:} Any sections/subsections marked with an asterisk are extra sections and have been added to this document as extensions of introduced theory.
\pagebreak

% \section{Determinants}

\section{Canonical Forms and Characteristic Values}
\subsection{Linear Transformations and Coordinates}
We will first recall linear transformations as homomorphisms between two vector spaces. We require the following foundational lemma:
\begin{lem}
An $n$-dimensional vector space $V$ is isomorphic to $\mathbb{F}^n$.
\end{lem}
\begin{proof}
Assume $\beta$ is an ordered basis for $V$. We claim that the coordinate function $[]_{\beta}$ is the linear isomorphism in question.

First, we must prove that it is a linear transformation, such that:
\begin{equation*}
[c_1v_1 + c_2v_2]_{\beta} = c_1[v_1]_{\beta} + c_2[v_2]_{\beta}
\end{equation*}
for $v_1, v_2 \in V$. Let $\beta = {b_1, \cdots b_n}$, then:
\begin{equation*}
v_1 = a_1b_1 + \cdots + a_nb_n \land v_2 = d_1b_1 + \cdots + d_nb_n
\end{equation*}
\begin{equation*}
\implies c_1v_1 + c_2v_2 = (c_1a_1 + c_2d_1)b_1 + \cdots + (c_1a_n + c_2d_n)b_n
\end{equation*}
\begin{equation*}
[c_1v_1 + c_2v_2]_{\beta} = \begin{bmatrix}
c_1a_1 + c_2d_1\\ \vdots \\ c_1a_n + c_2d_n
\end{bmatrix} = \begin{bmatrix}
c_1a_1\\ \vdots \\ c_1a_n
\end{bmatrix} + \begin{bmatrix}
c_2d_1 \\ \vdots \\ c_2d_n
\end{bmatrix} = c_1\begin{bmatrix}
a_1 \\ \vdots \\ a_n
\end{bmatrix} + c_2\begin{bmatrix}
d_1 \\ \vdots \\ d_n
\end{bmatrix} = c_1[v_1]_{\beta} + c_2[v_2]_\beta
\end{equation*}
To prove injectivity, assume $[v]_\beta = [w]_\beta$, then:
\begin{equation*}
[v]_\beta - [w]_\beta = \begin{bmatrix}
a_1 - d_1\\ \vdots \\ a_n - d_n
\end{bmatrix} = \begin{bmatrix}
0 \\ \vdots \\ 0
\end{bmatrix}
\end{equation*}
Thus for each $i \in {1, \cdots, n}$, we have that $a_i = d_i$, so by the uniqueness of linear combinations for independent vectors, $v = w$.

Finally, surjectivity follows from the fact that:
\begin{equation*}
\forall x \in \mathbb{F}^n: x_1b_1 + \cdots + x_nb_n \in V
\end{equation*}
by closure under addition and scalar multiplication. Thus $[x_1b_1 + \cdots + x_nb_n]_\beta = x$.
\end{proof}
Now consider the linear transformation:
\begin{equation*}
T: V \xrightarrow{} W
\end{equation*}
with the ordered basis $\beta_V$ and $\beta_W$. Then for some $v \in V$:
\begin{equation*}
v = c_1v_1 + \cdots c_nv_n \implies Tv = c_1Tv_1 + \cdots + c_nTv_n
\end{equation*}
Thus $[v]_{\beta_V} = [c_1, \cdots, c_n]^T$ and furthermore:
\begin{equation*}
[Tv]_{\beta_W} = [c_1Tv_1 + \cdots + c_nTv_n]_{\beta_W} = c_1[Tv_1]_{\beta_W} + \cdots + c_n[Tv_n]_{\beta_W}
\end{equation*}
Thus, we can represent $T$ as the matrix:
\begin{equation*}
[T]_{\beta_V}^{\beta_W} \coloneq \begin{bmatrix}
[Tv_1]_{\beta_W} & \cdots & [Tv_n]_{\beta_W}
\end{bmatrix}
\end{equation*}
So, we see that:
\begin{equation*}
[Tv]_{\beta_W} = [T]_{\beta_V}^{\beta_W}[v]_{\beta_V}
\end{equation*}

We now explore transformation composition and its interaction with choosing bases. Consider:
\begin{equation*}
T: V \xrightarrow{} W \land Q: W \xrightarrow{} U
\end{equation*}
with $\beta_V, \beta_W, \beta_U$ being ordered bases for their respective vector spaces. We know by the previous result:
\begin{equation*}
[T]_{\beta_V}^{\beta_W}[v]_{\beta_V} = [Tv]_{\beta_W}
\end{equation*}
Then, adding on $Q$, we have that:
\begin{equation*}
[Q]_{\beta_W}^{\beta_U}[Tv]_{\beta_W} = [QTv]_{\beta_U} = [QT]_{\beta_V}^{\beta_U}[v]_{\beta_V}
\end{equation*}
\begin{equation*}
\implies [Q]_{\beta_W}^{\beta_V} \circ [T]_{\beta_V}^{\beta_W} = [QT]_{\beta_V}^{\beta_U}
\end{equation*}
We can essentially write this as a commutative diagram. The result allows us to introduce the following:

\begin{defn}[Similar Matrices]
Two matrices are similar or $A \sim B$ if there exists a matrix $P$ such that:
\begin{equation*}
A = PBP^{-1}
\end{equation*}
\end{defn}
\begin{lem}
Given a linear operator $T: V \xrightarrow{} V$ (or even a linear transform $T: V \xrightarrow{} W$ where $\dim(V) = \dim(W)$), selecting two ordered bases $\alpha, \beta$, we have that:
\begin{equation*}
[T]_{\alpha}^{\alpha} \sim [T]_\beta^\beta
\end{equation*}
\end{lem}
\begin{proof}
The logic comes from the idea of change of bases and linear operators. Consider the following commutative diagram and the ordered bases for $\beta, \alpha$ for $V$ with $T: V \xrightarrow{} V$:
\begin{center}
\begin{tikzcd}
v \arrow[r, "I"] \arrow[d, "\lbrack \rbrack_\alpha"'] & v  \arrow[r, "T"] \arrow[d, "\lbrack \rbrack_\beta"'] & Tv \arrow[d, "\lbrack \rbrack_\beta"] \arrow[r, "I"] & Tv \arrow[d, "\lbrack \rbrack_\alpha"'] \\
\lbrack v \rbrack_\alpha \arrow[r, "\lbrack I \rbrack_{\alpha}^{\beta}"] & \lbrack v \rbrack_\beta \arrow[r, "\lbrack T \rbrack_{\beta}^{\beta}"]   & \lbrack Tv \rbrack_\beta \arrow[r, "\lbrack I \rbrack_{\beta}^{\alpha}"] & \lbrack Tv \rbrack_\alpha
\end{tikzcd}
\end{center}
Given the fact that $[T]_{\alpha}^\alpha [v]_\alpha = [Tv]_\alpha$, we see that by the commutative diagram:
\begin{equation*}
[T]_{\alpha}^{\alpha} = [I]_{\beta}^\alpha \circ [T]_{\beta}^{\beta} \circ [I]_{\alpha}^{\beta}
\end{equation*}
We observe the following, applying our results on the composition of multiple linear transforms:
\begin{equation*}
[I]_{\beta}^\alpha \circ [I]_\alpha^\beta = [I]_\alpha^\alpha = I
\end{equation*}
\begin{equation*}
[I]_\alpha^\beta \circ [I]_\beta^\alpha = [I]_\beta^\beta = I
\end{equation*}
Thus we can call $P \coloneq [I]_\beta^\alpha$, and we get that $P^{-1} = [I]_\alpha^\beta$. Then:
\begin{equation*}
A \coloneq [T]_\alpha^\alpha \land B \coloneq [T]_\beta^\beta
\end{equation*}
\begin{equation*}
\implies A = PBP^{-1}
\end{equation*}
\end{proof}
 
\subsection{Characteristic Values}
For this set of notes, we shall continue to use the terminology of "characteristic values" which readers must note is synonymous to eigenvalues.

\begin{defn}
A linear operator $T$ is simply a linear map from a vector space to itself:
\begin{equation*}
T: V \xrightarrow{} V
\end{equation*}
We say that $v$ is a characteristic vector with an associated characteristic value $c \in F$ if:
\begin{equation*}
Tv = cv
\end{equation*}
The set of all vectors that follow this property for a certain $c$ is called the characteristic space of $c$ and is equivalently $\ker(cI - T)$.
\end{defn}
We already see that if some $v$ is in the characteristic space of $c$, then its scalar multiples are too:
\begin{equation*}
\forall b \in F: T(bv) = b(Tv) = b(cv) = c(bv)
\end{equation*}
We also see that the addition of two vectors in the characteristic space of $c$ is also in the characteristic space:
\begin{equation*}
T(v + w) = T(v) + T(w) = cv + cw = c(v+w)
\end{equation*}
The $0$ vector is trivially in the characteristic space of any arbitrary $c$. Thus the characteristic space of $c$ is a subspace of $V$. If the characteristic space of $c$ is the zero subspace, then it is not a characteristic value.

We will denote the characteristic space of $c$ as $E_c$ from now on.

\begin{lem}
If $0$ is a characteristic value of a linear operator $T$, then it has a non-trivial kernel.
\end{lem}
\begin{proof}
This follows trivially as this means there exists some subspace $E_0$ (which is not the zero subspace), such that:
\begin{equation*}
\forall v \in E_0: Tv = 0v = 0 \implies E_0 \subseteq \ker(T)
\end{equation*}
\end{proof}

\begin{lem}
Given a linear transformation $T: V \xrightarrow{} W$:
\begin{equation*}
v \in \ker(T)
\end{equation*}
if and only if:
\begin{equation*}
[v]_{\beta_V} \in \ker([T]_{\beta_V}^{\beta_W})
\end{equation*}
and $\dim(\ker(T)) = \dim(\ker([T]_{\beta_V}^{\beta_W}))$.
\end{lem}
\begin{proof}
$\implies:$ If $v \in \ker(T)$, then $Tv = 0$, which tells us that:
\begin{equation*}
[Tv]_{\beta_W} = [T]_{\beta_V}^{\beta_W}[v]_{\beta_V} = [0]_{\beta_W}
\end{equation*}

$\impliedby:$ If $[v]_{\beta_V} \in \ker([T]_{\beta_V}^{\beta_W})$, then:
\begin{equation*}
[T]_{\beta_V}^{\beta_W}[v]_{\beta_V} = [Tv]_{\beta_W} = [0]_{\beta_W}
\end{equation*}
which by injectivity of the coordinate function implies that $Tv = 0$, so $v \in \ker(T)$. Now let $\mathcal{K}$ be a basis for $\ker(T)$ with vectors $k_1, \cdots, k_m$, then we claim that:
\begin{equation*}
\mathcal{K}' \coloneq \{[k_1]_{\beta_V}, \cdots [k_m]_{\beta_V}\}
\end{equation*}
is a basis for $\ker([T]_{\beta_V}^{\beta_W})$:
\begin{equation*}
\forall [v]_{\beta_V} \in \ker([T]_{\beta_V}^{\beta_W}): [Tv]_{\beta_W} = [0]_{\beta_W}
\end{equation*}
tells us that once again $v \in \ker(T)$, so it can be represented as a linear combination of $\mathcal{K}$ such that:
\begin{equation*}
\implies [T]_{\beta_V}^{\beta_W}(c_1[k_1]_{\beta_V} + \cdots + c_m[k_m]_{\beta_V}) = [0]_{\beta_W}
\end{equation*}
We also note that $\mathcal{K}'$ is linearly independent:
\begin{equation*}
c_1[k_1]_{\beta_V} + \cdots + c_m[k_m]_{\beta_V} = [0]_{\beta_V} \implies [c_1k_1 + \cdots + c_mk_m]_{\beta_V} = [0]_{\beta_V}
\end{equation*}
And by injectivity of the coordinate function, since $\mathcal{K}$ is linearly independent so that:
\begin{equation*}
\implies c_1k_1 + \cdots + c_mk_m = 0 \implies c_1, \cdots c_m = 0
\end{equation*}
Thus $\mathcal{K}'$ is the basis we are looking for, showing the dimensions are equal. 
\end{proof}

\begin{thm}
The following are equivalent conditions for any arbitrary basis $\beta$ of the vector space $V$:
\begin{enumerate}
\item $c$ is a characteristic value of $T$
\item $cI - T$ is noninvertible
\item $\det(cI - [T]_{\beta}^{\beta}) = 0$
\end{enumerate}
\end{thm}
\begin{proof}
$(1 \implies 2):$ If $c$ is a characteristic value, then we have a subspace $E_c$ such that:
\begin{equation*}
\forall v \in E_c: Tv  = cv \implies (cI - T)v = 0
\end{equation*}
Because $cI - T$ has a nontrivial kernel, it is not invertible (since it is not an injection).

$(2 \implies 3):$ This follows from determinant properties.

$(3 \implies 1):$ If $\det(cI - [T]_\beta^\beta) = 0$, then the kernel is nontrivial, so:
\begin{equation*}
\exists v \in \ker(T): (cI - [T]_\beta^\beta)[v]_\beta = 0 \implies [Tv]_\beta = [cv]_\beta
\end{equation*}
By the injectivity proved in Lemma 1.1, we obtain:
\begin{equation*}
Tv = cv
\end{equation*}
which tells us by Definition 2.2, that $c$ is a characteristic value of the operator. 
\end{proof}

In order to find the characteristic values, we can use the characteristic polynomial:
\begin{defn}
The characteristic polynomial for a linear operator $T$ is defined as $q_T(x) \coloneq \det(xI - [T]_\beta^\beta)$ (or alternatively $\det(T - xI)$). The roots to such a polynomial are equivalent to the characteristic values
by Theorem 1.1.
\end{defn}

Using the theory of similar matrices and determinants, we can also obtain the following result:
\begin{lem}
The characteristic polynomial is independent of the chosen basis $\beta$ of $V$, despite a different matrix representation.
\end{lem}
\begin{proof}
Consider two ordered bases for $V$, which we call $\alpha, \beta$. Then $q_T(x) \coloneq \det(xI - [T]_\beta^\beta)$.

By Lemma 2.2, we see that there exists some square matrix $P$ such that:
\begin{equation*}
q_T'(x) \coloneq \det(xI - [T]_\alpha^\alpha) = \det(xI - P[T]_\beta^\beta P^{-1}) = \det(xIPP^{-1} - P[T]_\beta^\beta P^{-1})
\end{equation*}
\begin{equation*}
= \det((xIP - P[T]_\beta^\beta)P^{-1}) = \det((xPI - P[T]_\beta^\beta)P^{-1}) = \det(PxI - P[T]_\beta^\beta)\det(P^{-1})
\end{equation*}
\begin{equation*}
= \det(P)\det(xI - [T]_\beta^\beta)\det(P^{-1}) = \det(xI - [T]_\beta^\beta) = q_T(x)
\end{equation*}
Thus the characteristic polynomial is the same.
\end{proof}

\begin{lem}
The degree of the characteristic polynomial is $\dim(V)$.
\end{lem}
\begin{proof}
This comes from determinant properties, whereby for some $n \times n$ square matrix $A$:
\begin{equation*}
\det(A) \coloneq \sum_{\sigma \in S_n}\text{sgn}(\sigma)\prod_{i=1}^{n}A_{i, \sigma(i)}
\end{equation*}
where $S_n$ is the symmetric group whose elements are permutations $\sigma$. If we consider the identity permutation, then we obtain the product:
\begin{equation*}
\prod_{i=1}^n(xI - [T]_{\beta}^\beta)_{i, i} = \prod_{i=1}^n(x - ([T]_{\beta}^{\beta})_{i, i})
\end{equation*}
which is the highest degree term of the characteristic polynomial which is $\dim(V)$ as $[T]_\beta^\beta \in \mathcal{M}_{\dim(V) \times \dim(V)}(F)$.
\end{proof}

We say that a linear operator $T$ is diagonalizable if we can obtain a set of characteristic vectors from the characteristic spaces forms a basis for $V$:
\begin{equation*}
\beta \coloneq \{v_1, \cdots, v_n\}
\end{equation*}
where $Tv_i = c_iv_i$, telling us that $[Tv_i]_\beta = [c_iv_i]_\beta = c_i[v_i]_\beta = c_ie_i$ by Lemma 2.1:
\begin{equation*}
[T]_\beta^\beta = \begin{bmatrix}
c_1 & \cdots & 0\\
\vdots & \ddots & \vdots\\
0 & \cdots & c_n
\end{bmatrix}
\end{equation*}

\begin{defn}[Multiplicities]
The algebraic multiplicity of a characteristic value is the number of times it appears as the root of the characteristic polynomial $q_T(x)$.

The geometric multiplicity of a characteristic value is the dimension of its characteristic space $E_c$.
\end{defn}

\begin{lem}
The geometric multiplicity of a characteristic value $c$ is always less than or equal to the algebraic multiplicity.
\end{lem}
\begin{proof}
We defined the geometric multiplicity as $g = \dim(\ker(cI - T))$. Since $\ker(cI - T)$ is a subspace of $V$, we define a basis for it as $\{b_1, ..., b_g\}$. Then, we extend this basis to get:
\begin{equation*}
\beta_V \coloneq \{b_1, ..., b_g, b_{g+1}, ..., b_n\}
\end{equation*}
which is a basis for $V$. Now observe that:
\begin{equation*}
[T]_{\beta_V}^{\beta_V} = \begin{bmatrix}
cI_g & C\\
0 & D
\end{bmatrix}
\end{equation*}
where $I_g$ is the $g \times g$ identity, $0$ is the $(n-g)\times g$ matrix with entries of $0$. $C$ is a $g \times (n-g)$ matrix and $D$ is the remaining $(n-g) \times (n-g)$ matrix.

\begin{equation*}
q_T(x) = \det(xI - [T]_{\beta_V}^{\beta_V}) = (x-c)^g\det(xI_{n-g} - D)
\end{equation*}
Thus $(x-c)^g$ divides the characteristic polynomial, so we have that algebraic multiplicity is greater than or equal to $g$ (depending on how the remainder determinant term factors). By definition of characteristic value and geometric multiplicity in Theorem 1.1, we have that:
\begin{equation*}
1 \leqslant \dim(ker(cI - T))
\end{equation*}
\end{proof}

\begin{defn}
Given two vector subspaces $V, W \subseteq U$, we define their sum as:
\begin{equation*}
V + W \coloneq \{v + w : v \in V \land w \in W\}
\end{equation*}
We can write this as a direct sum:
\begin{equation*}
V + W = V \oplus W
\end{equation*}
if and only if:
\begin{equation*}
\forall v \in V, \forall w \in W: v+w = 0 \implies v, w = 0
\end{equation*}
We call this linear independence of vector subspaces, which can be extended to finite sums of vector subspaces.
\end{defn}
We prove a small equivalence of conditions for the two subspace condition, allowing us to expand to the $k$-subspace condition in Lemma 1.8.

We see that if $0 \subset V \cap W$, then there exists some $v \neq 0$ such that:
\begin{equation*}
v \in V \land v \in W
\end{equation*}
So take $-v \in V, v \in W$, then $-v, v \neq 0$ but $-v + v = 0$ proving that linear independence implies:
\begin{equation*}
V \cap W = 0
\end{equation*}
Now for the converse, if there exists some $v \in V, w \in W$ such that $v+w = 0 \land v, w \neq 0$, then $v = -w$.

Since $-w \in V$, then we have that $w \in V \land w \in W$, so $0 \subset V \cap W$. We will also require the fact that the sum of subspaces is a subspace, which we can prove easily.

Given $W_1, ..., W_k$ subspaces, then $\forall j \in 1, ..., k: 0 \in W_j$, thus:
\begin{equation*}
0 \in W_1 + \cdots + W_k
\end{equation*}
Now take some $u, v \in W_1 + \cdots + W_k$. We see that by definition of subspace addition:
\begin{equation*}
\exists w_1, ..., w_k: u = w_1 + \cdots + w_k
\end{equation*}
\begin{equation*}
\exists w_1', ..., w_k': v = w_1' + \cdots w_k'
\end{equation*}
\begin{equation*}
\implies u + v = (w_1 + w_1') + \cdots + (w_k + w_k') \in W_1 + \cdots + W_k
\end{equation*}
by subspace proeprty of each $W_j$. Scalar multiplication follows similarly. We also see that under a direct sum, the representation of a vector $v$ must be a unique combination by linear independence of subspaces.
\begin{lem}
The following statements are equivalent regarding direct sums, given subspaces $W_1, \cdots W_k$:
\begin{enumerate}
\item $W_1, \cdots, W_k$ are linearly independent
\item The following condition:
\begin{equation*}
\forall j \in {1, \cdots, k}: W_j \cap \sum_{i \neq j}W_i = 0
\end{equation*}
\item Call the bases for each subspace $\beta_1, \cdots, \beta_k$, then for:
\begin{equation*}
W = W_1 + \cdots + W_k
\end{equation*}
\begin{equation*}
\beta_W = \beta_1 \cup \cdots \cup \beta_k
\end{equation*}
is a valid basis for $W$.
\end{enumerate}
\end{lem}
\begin{proof}
$(1 \implies 2)$: If the subspaces are linearly independent, then:
\begin{equation*}
\forall w_i \in W_i: w_1 + \cdots + w_k = 0 \implies w_1, \cdots, w_k = 0
\end{equation*}
Assume there exists some $j \in 1, ..., k$ such that:
\begin{equation*}
0 \subset W_j \cap \sum_{i \neq j}W_i
\end{equation*}
Then we can say there is some $v \neq 0$ such that:
\begin{equation*}
v \in W_j \cap \sum_{i \neq j}W_i \implies v \in W_j \land -v \in W_j
\end{equation*}
It also tells us that $v$ is in the sum of the $W_i$ subspaces excluding $j$, so there are some vectors $w_i$ such that:
\begin{equation*}
v = \sum_{i \neq j}w_i \land w_j = -v \implies \exists w_i \in W_i: w_j + \sum_{i \neq j}w_i = -v + v = 0
\end{equation*}
yet not all of the $w_i$ vectors are the $0$ vector.

$(2 \implies 3)$: Take some $w \in W$, then there exists $w_i \in W_i$ for $i \in 1, ..., k$ such that:
\begin{equation*}
w = w_1 + \cdots + w_k = (c_1^{(1)}b_1^{(1)} + \cdots + c_{m(1)}^{(1)}b_{m(1)}^{(1)}) + \cdots + (c_1^{(k)}b_1^{(k)} + \cdots + c_{m(k)}^{(k)}b_{m(k)}^{(k)})
\end{equation*}
where $\dim(W_j) = m(j)$ and the superscripts identifying which subspace's bases and coordinates the vectors and scalars represent respectively.

We see that each $w \in W$ can be represented as a linear combination of:
\begin{equation*}
\beta_1 \cup \cdots \cup \beta_k
\end{equation*}
We now aim to show the basis vectors are linearly independent to each other. If $w$ were the $0$ vector in this case, so that:
\begin{equation*}
(c_1^{(1)}b_1^{(1)} + \cdots + c_{m(1)}^{(1)}b_{m(1)}^{(1)}) + \cdots + (c_1^{(k)}b_1^{(k)} + \cdots + c_{m(k)}^{(k)}b_{m(k)}^{(k)}) = 0
\end{equation*}
for arbitrary constants in $F$, then we can prove this iteratively for all $j \in 1, ..., k$. Consider $j = 1$, then using the results from our mini proof after Definition 1.5:
\begin{equation*}
W_1 \cap \sum_{i \neq 1}W_i = 0 \implies (c_1^{(1)}b_1^{(1)} + \cdots + c_{m(1)}^{(1)}b_{m(1)}^{(1)}) = 0
\end{equation*}
and
\begin{equation*}
(c_1^{(2)}b_1^{(2)} + \cdots + c_{m(2)}^{(2)}b_{m(2)}^{(2)}) + \cdots + (c_1^{(k)}b_1^{(k)} + \cdots + c_{m(k)}^{(k)}b_{m(k)}^{(k)}) = 0
\end{equation*}
By linear independence of $\beta_1$, we have that the first set of constants are equal to $0$. We can repeatedly iterate this for all $j \in 2, ..., k$ and get that the union of bases is a basis for the sum.

$(3 \implies 1)$: Assume $\beta_1 \cup \cdots \cup \beta_k$ is a valid basis for $W$, then:
\begin{equation*}
\forall w_1, ..., w_k: w_1 + \cdots + w_k = 0 \implies (c_1^{(1)}b_1^{(1)} + \cdots + c_{m(1)}^{(1)}b_{m(1)}^{(1)}) + \cdots + (c_1^{(k)}b_1^{(k)} + \cdots + c_{m(k)}^{(k)}b_{m(k)}^{(k)}) = 0
\end{equation*}
Because the union of bases is a basis for $W$, then all the constants are $0$ by linear independence, thus all $w_i = 0$, proving linear independence of the subspaces.
\end{proof}
The third equivalent condition will be important for the next theorem. A useful consequence of the third equivalent notion is that:
\begin{equation*}
\dim(W) = \dim(W_1) + \cdots + \dim(W_k)
\end{equation*}
which comes from the fact that the union of bases is a basis for the sum.

\begin{thm}
The following statements are equivalent to diagonalizability:
\begin{enumerate}
\item The linear operator $T: V \xrightarrow{} V$ is diagonalizable
\item If $d_1, \cdots d_k$ are the algebraic multiplicities, then:
\begin{equation*}
q_T(x) = (x - c_1)^{d_1} \cdots (x - c_k)^{d_k}
\end{equation*}
and $d_i = \dim(E_{c_i})$ (algebraic multiplicity is equal to geometric multiplicity)
% \item The sum of geometric multiplicities adds up to the dimension of the vector space:
% \begin{equation*}
% \dim(V) = \sum_{i=1}^{k}\dim(E_{c_i})
% \end{equation*}
\end{enumerate}
\end{thm}
\begin{proof}
$(\implies)$: Assume $d_1, \cdots, d_k$ are the algebraic multiplicities of the characteristic values (where $k \leqslant n$), then given a basis $\beta$ composed of characteristic vectors, if we order the vectors in a certain way, we get:
\begin{equation*}
[T]_\beta^\beta = \begin{bmatrix}
c_1I_1 & \cdots & 0\\
\vdots & \ddots & \vdots\\
0 & \cdots & c_kI_k
\end{bmatrix}
\end{equation*}
where $I_j$ is a $d_j \times d_j$ matrix. If $k = n$, then each $d_i = 1$ and we have $c_n$ characteristic values. If $k < n$, then some of the characteristic values must be shared by some of the vectors, so we can order the basis accordingly.

Then:
\begin{equation*}
\det(xI - [T]_\beta^\beta) = (x - c_1)^{d_1} \cdots (x - c_k)^{d_k}
\end{equation*}
We now must prove that the algebraic multiplicit of each characteristic value is equal to its geometric multiplicity, which we do via Definition 1.2 and Lemma 1.4:
\begin{equation*}
E_{c_i} \coloneq \ker(c_iI - T) \implies \dim(E_{c_i}) = \dim(\ker(c_iI - T)) = \dim(\ker(c_iI - [T]_{\beta}^{\beta}))
\end{equation*}
Observe that $c_iI - [T]_{\beta}^{\beta}$ (by the uniqueness of the characteristic values) only has $d_i$ colums of $0$ vectors with the rest being linearly independent, so:
\begin{equation*}
\dim(\ker(c_iI - [T]_{\beta}^\beta)) = \dim(E_{c_i}) = d_i
\end{equation*}

$(\impliedby):$ For this proof, we first need to prove:
\begin{equation*}
\dim(V) = \sum_{i=1}^{k}d_i
\end{equation*}
This follows from Lemma 1.6. We have that $\dim(V) = \deg(q_T(x)) = \sum_{i=1}^{k}d_i$. Using the second part of statement $2$, we obtain:
\begin{equation*}
\dim(V) = \sum_{i=1}^{k}\dim(E_{c_i})
\end{equation*}

% $(3 \implies 1)$: 
We aim to show that:
\begin{equation*}
E_{c_1} + \cdots + E_{c_k} = E_{c_1} \oplus \cdots \oplus E_{c_k}
\end{equation*}
We show it via condition 1 in Lemma 1.8. We want to show that:
\begin{equation*}
\forall j \in 1, ..., k: \sum_{i=1}^k v_j = 0 \implies v_j = 0
\end{equation*}
Define:
\begin{equation*}
p(x) \coloneq \prod_{i=1}^{k}(c_i - x) \implies p(T) = \prod_{i=1}^{k}(c_iI - T)
\end{equation*}
In order to proceed with this proof, we need to prove that the factors of $p(T)$ commute. Take an arbitrary $c_i, c_j$, then:
\begin{equation*}
(c_iI - T)(c_jI - T) = c_ic_jI - c_iT - c_jT + T^2 = c_jc_iI - c_jT - c_iT + T^2 = (c_jI - T)(c_iI - T)
\end{equation*}
Observe that we used the commutativity of multiplication in a field $F$. If this were a module over a ring $R$, then the factors of $p(T)$ may not commute.

We define:
\begin{equation*}
\forall j \in 1, ..., k: p_j(T) \coloneq \prod_{i \neq j}(c_iI - T)
\end{equation*}
So, by commuting factors accordingly, we can eliminate all $v_i$ vectors except $v_j$:
\begin{equation*}
\sum_{i \neq j}p_j(T)v_i = p(T)0 = 0 \implies p_j(T)v_j = 0
\end{equation*}
We apply $v_j$ to each and obtain:
\begin{equation*}
\prod_{i \neq j}(c_i - c_j)v_j = 0 \implies v_j = 0
\end{equation*}
We can do this for all $j \in 1, ..., k$, implying that all $v_j$ must be zero, showing the characteristic spaces are linearly independent.

As a result, we have a basis for $E$ such that:
\begin{equation*}
\beta_E = \beta_1 \cup \cdots \cup \beta_k
\end{equation*}
where $E$ is a subspace of $V$. We also have the consequence that:
\begin{equation*}
\dim(E) = \dim(E_{c_1}) + \cdots + \dim(E_{c_k}) = \dim(V)
\end{equation*}
Since the dimensions are equal, that means $\beta_E$ being $\dim(V)$ linearly independent vectors means they span $V$ and are a valid basis for $V$.

So, we see that $[T]_{\beta_E}^{\beta_E}$ is a basis that diagonalizes the matrix as:
\begin{equation*}
\forall j \in 1, ..., k, \forall b_i^{(j)} \in \beta_E: Tb_i^{(j)} = c_jb_i^{(j)}
\end{equation*}
or each of the elements in $\beta_E$ come from characteristic spaces.
\end{proof}
The important takeaways are that a linear operator is diagonalizable if and only if the algebraic multiplicity of each characteristic value is equal to its geometric multiplicity, which requires the fact that
the characteristic spaces are linearly independent (or to put informally, the characteristic spaces have no overlap).

By matrix similarity in Lemma 1.2, if $T$ is diagonalizable such that:
\begin{equation*}
[T]_\beta^\beta
\end{equation*}
is diagonal and $\beta$ is a basis of characteristic vectors spanning $V$. Then taking any arbitrary basis $\alpha$ of $V$, we see that:
\begin{equation*}
[T]_\beta^\beta = [I]_\alpha^\beta \circ [T]_\alpha^\alpha \circ [I]_\beta^\alpha
\end{equation*}
We call $[I]_\beta^\alpha$ (or the matrix of characteristic basis vectors represented in terms of $\alpha$) the modal matrix $M$.

A useful result is that the powers of similar matrices are also similar. So, if a linear operator is diagonalizable, then:
\begin{equation*}
[T]_\alpha^\alpha = [I]_\beta^\alpha \circ [T]_\beta^\beta \circ [I]_\alpha^\beta
\end{equation*}
\begin{equation*}
\implies ([T]_\alpha^\alpha)^n = [I]_\beta^\alpha \circ ([T]_\beta^\beta)^n \circ [I]_\alpha^\beta
\end{equation*}
Since $\beta$ is a characteristic basis, the power of a diagonal matrix is simply its terms on the diagonal raised to that power.

\subsubsection{Dynamical Systems Problem}
We examine state space representations, which are key to modern control.

\begin{defn}
The state variables are the smallest set of linearly independent variables that describe a system's state.
\end{defn}
While this is quite vague of a definition, we can be more specific saying that typically state variables are the configuration variables plus velocities. Just like the configuration space, the state space can be a manifold as well.

For the purpose of these notes, we shall consider the state space $\mathcal{S}$ to be a euclidean space $\mathbb{R}^n$, so that states are vectors. The additional condition is that because states are evolving, the states must have an initial condition, which combined with the update rule provides us a way to calculate the most recent state.
We can define the control space $\mathcal{U}$ in a similar fashion.

Note that $x(t)$ is a function which takes in time and outputs our state, so we are interested in solving for $x(t)$ to understand how the system's state progresses based on some applied control and time.

We also have a control input defined for all time, which we define $u(t): \mathbb{R}_{\geqslant 0} \xrightarrow{} \mathbb{R}^m$ where:
\begin{equation*}
u(t) \coloneq \begin{bmatrix}
u_1(t)\\
\vdots\\
u_m(t)
\end{bmatrix}, \ x(t) \coloneq \begin{bmatrix}
x_1(t)\\
\vdots\\
x_n(t)
\end{bmatrix}
\end{equation*}
such that $\forall j \in 1, ..., m: u_j: \mathbb{R}_{\geqslant 0} \xrightarrow{} \mathbb{R}$ and $x_1, ..., x_n$ are the state variables. Then we claim there exists a vector valued function in $C^1$ on its domain:
\begin{equation*}
f: \mathbb{R}^{n+m+1} \xrightarrow{} \mathbb{R}^n
\end{equation*}
such that:
\begin{equation*}
\forall t \in \mathbb{R}_{\geqslant 0}: x'(t) = f(x(t), u(t), t) = \begin{bmatrix}
f_1(x(t), u(t), t)\\
\vdots\\
f_n(x(t), u(t), t)
\end{bmatrix}
\end{equation*}
such that $\forall i \in 1, ..., n$, $f_i$ is a scalar valued function (this takes the form as showed in the differential equations notes).

For now we will eliminate the extra $t$ variable in the domain, so that we have $f: \mathbb{R}^{n+m} \xrightarrow{} \mathbb{R}^n$. Here we can make both algebraic and analytical connections:
\begin{lem}
If $f$ is a linear transformation, then the dynamical system can be represented as:
\begin{equation*}
x'(t) = Ax(t) + Bu(t)
\end{equation*}
where $A \in \mathcal{M}_{n \times n}(\mathbb{R})$ and $B \in \mathcal{M}_{n \times m}(\mathbb{R})$ and an associated initial condition.
\end{lem}
\begin{proof}
If $f$ is a linear transformation, then (considering with respect to the canonical bases), it admits a matrix representation:
\begin{equation*}
x'(t) = \begin{bmatrix}
f(e_1) & \cdots & f(e_{n+m})
\end{bmatrix}\begin{bmatrix}
x(t)\\u(t)
\end{bmatrix} = \begin{bmatrix}
f(e_1) & \cdots & f(e_{n})\end{bmatrix}x(t) + \begin{bmatrix}f(e_{n+1}) & \cdots & f(e_{n+m})\end{bmatrix}u(t)
\end{equation*}
So we set $A, B$ accordingly and we get the desired form. Observe that if we select a different basis for the state space or even the control space, then the coordinate representation of the states and inputs change. 
\end{proof}
Such a system is called a continuous time invariant system. Consider the simple case:
\begin{lem}
Consider a CTI system with no applied control input such that:
\begin{equation*}
x'(t) = Ax(t), \ x(0) = x_0
\end{equation*}
with $A$ diagonalizable. Then the limit of the state as $t \xrightarrow{} \infty$ exists if all characteristic values are less than or equal to $0$.
\end{lem}
\begin{proof}
By definition of diagonalizable, we have a characteristic basis $\beta \coloneq \{b_1, \cdots, b_n\}$ that spans the state space. Additionally:
\begin{equation*}
\begin{bmatrix}
x_1'(t)\\
\vdots\\
x_n'(t)  
\end{bmatrix} = A\begin{bmatrix}
x_1(t)\\
\vdots\\
x_n(t)
\end{bmatrix}
\end{equation*}
The family of solutions based on the initial condition is:
\begin{equation*}
S \coloneq \left\{\sum_{i=1}^{n}d_ib_ie^{c_it} : \sum_{i=1}^{n}d_ib_i = x_0\right\}
\end{equation*}
so for $\lim_{t \xrightarrow{} \infty}x(t)$ to exist, it must be that:
\begin{equation*}
\forall i \in 1, ..., n: c_i \leqslant 0
\end{equation*}
but $0$ cannot be a characteristic value, so it means $c_i < 0$. 
\end{proof}
We now consider a discrete time invariant version of this relationship with no applied control input:
\begin{lem}
Consider a DTI system described by the recursive time-step relation:
\begin{equation*}
x_{t+1} = Ax_t, \ x(0) = x_0
\end{equation*}
with $A$ once again diagonalizable. Then the limit of the state as $t \xrightarrow{} \infty$ exists if the characteristic values are in the interval $(-1, 1]$. 
\end{lem}
\begin{proof}
We can iterate the relation and obtain the particular solution:
\begin{equation*}
x_t = A^tx_0
\end{equation*}
Since $A$ is diagonalizable, there exists a characteristic basis $\beta$ spanning the state space so that:
\begin{equation*}
x_t = A^t(d_1b_1 + \cdots + d_nb_n) = d_1A^tb_1 + \cdots d_nA^tb_n = d_1c_1^tb_1 + \cdots + d_nc_n^tb_n
\end{equation*}
where the $c_i$ are the characteristic values. Then:
\begin{equation*}
\lim_{t \xrightarrow{} \infty}x_t = \sum_{i=1}^{n}d_i \cdot \lim_{t \xrightarrow{} \infty}c_i^tb_i
\end{equation*}
exists if for all characteristic values are between $(-1, 1]$.
\end{proof}
\end{document}